% Mechanistic Validity — NEMI 2026 poster, v6
% Five rows, one per validity type.
% Left: 3-column criteria table (ID, definition, source).
% Right: two case study examples.
%
%   cd poster && lualatex mechval_poster_v6.tex

\documentclass[25pt, a0paper, portrait]{tikzposter}

\usepackage[T1]{fontenc}
\usepackage{amsmath,amssymb}
\usepackage{array}
\usepackage{booktabs}
\usepackage{enumitem}
\usepackage{microtype}

\tikzposterlatexaffectionproofoff

% ── Colors ────────────────────────────────────────────
\definecolor{darktext}{HTML}{2C3E50}
\definecolor{blockbg}{HTML}{FFFFFF}
\definecolor{constructgreen}{HTML}{27AE60}
\definecolor{measureblue}{HTML}{2980B9}
\definecolor{internalorange}{HTML}{E67E22}
\definecolor{externalpurple}{HTML}{8E44AD}
\definecolor{interpteal}{HTML}{16A085}

% ── Theme ─────────────────────────────────────────────
\usetheme{Default}

\definecolorstyle{MechValColors}{
  \definecolor{colorOne}{HTML}{2C3E50}
  \definecolor{colorTwo}{HTML}{3498DB}
  \definecolor{colorThree}{HTML}{ECF0F1}
}{
  \colorlet{backgroundcolor}{colorThree}
  \colorlet{framecolor}{colorOne}
  \colorlet{titlefgcolor}{white}
  \colorlet{titlebgcolor}{colorOne}
  \colorlet{blocktitlebgcolor}{colorTwo}
  \colorlet{blocktitlefgcolor}{white}
  \colorlet{blockbodybgcolor}{blockbg}
  \colorlet{blockbodyfgcolor}{darktext}
  \colorlet{notefgcolor}{darktext}
  \colorlet{notebgcolor}{colorTwo!12}
  \colorlet{notefrcolor}{colorTwo}
}
\usecolorstyle{MechValColors}

\defineblockstyle{MechValBlock}{
  titlewidthscale=1.0, bodywidthscale=1.0, titleleft,
  titleoffsetx=0pt, titleoffsety=0pt,
  bodyoffsetx=0pt, bodyoffsety=0pt, bodyverticalshift=0pt,
  roundedcorners=12, linewidth=2pt, titleinnersep=10pt, bodyinnersep=12pt
}{
  \draw[rounded corners=\blockroundedcorners, color=blocktitlebgcolor,
        fill=blockbodybgcolor, line width=\blocklinewidth]
    (blockbody.south west) rectangle (blocktitle.north east);
  \ifBlockHasTitle
    \draw[rounded corners=\blockroundedcorners,
          fill=blocktitlebgcolor, color=blocktitlebgcolor,
          line width=\blocklinewidth]
      (blocktitle.south west) rectangle (blocktitle.north east);
  \fi
}
\useblockstyle{MechValBlock}

\definetitlestyle{MechValTitle}{
  width=\paperwidth, roundedcorners=0, linewidth=0pt, innersep=20pt,
  titletotopverticalspace=12mm, titletoblockverticalspace=18mm
}{
  \draw[fill=titlebgcolor, color=titlebgcolor]
    (\titleposleft, \titleposbottom)
    rectangle (\titleposright, \titlepostop);
}
\usetitlestyle{MechValTitle}

\title{%
  \parbox{0.92\linewidth}{\centering\bfseries
    Mechanistic Validity\\[6pt]
    {\large A Validity Framework for Mechanistic Interpretability}}}
\author{Elliot Tower\quad\texttt{elliot@elliottower.ai}}
\institute{}

\begin{document}

\maketitle[width=\paperwidth]

% ════════════════════════════════════════════════════════
% ROW 1: CONSTRUCT VALIDITY
% ════════════════════════════════════════════════════════
\begin{columns}
\column{0.42}

{
\colorlet{blocktitlebgcolor}{constructgreen}
\block{Construct Validity --- Is the target concept well-defined?}{
  A claim is bounded by its weakest validity dimension.
  Five types form a dependency chain: construct validity
  is a ceiling on all others.

  \vspace{4mm}
  \renewcommand{\arraystretch}{1.25}
  {\small
  \begin{tabular}{@{}c l p{0.38\linewidth} l@{}}
    \toprule
    \textbf{ID} & \textbf{Name} & \textbf{Tests} & \textbf{Source} \\
    \midrule
    C1 & Falsifiability & Can the claim be refuted? & Phil.\ of sci.\ \\
    C2 & Structural plaus.\ & Physically possible in the arch.? & Phil.\ of sci.\ \\
    C3 & Convergent val.\ & Do independent methods agree? & Psychometrics \\
    C4 & Discriminant val.\ & Distinguishes from neighbors? & Psychometrics \\
    C5 & Nomological val.\ & Fits a broader theory? & Psychometrics \\
    C6 & Complementation & Are subdivisions distinct? & Genetics \\
    \bottomrule
  \end{tabular}}
}
}

\column{0.29}

{
\colorlet{blocktitlebgcolor}{constructgreen}
\block{C4: Knowledge Neurons}{
  Neurons identified by activation magnitude, but the
  criterion cannot separate factual knowledge from
  linguistic regularity.

  \vspace{3mm}
  Suppressing these neurons produces off-target effects
  \textbf{worse than random neurons}.

  \vspace{3mm}
  Similarly, gender bias circuits (Vig 2020) identify
  heads correlated with gendered outputs but never test
  whether those heads encode bias or linguistic competence.
}
}

\column{0.29}

{
\colorlet{blocktitlebgcolor}{constructgreen}
\block{C6: IOI Complementation}{
  IOI names ``name movers,'' ``S-inhibition heads,'' and
  ``backup name movers.'' Whether these are distinct
  functional units has never been tested.

  \vspace{3mm}
  Benzer's cis-trans test (1955): run two passes, one
  with each component intact, construct a hybrid. If
  the hybrid recovers, the components are independent.
  If it fails, they are the same functional unit.
}
}

\end{columns}

% ════════════════════════════════════════════════════════
% ROW 2: MEASUREMENT VALIDITY
% ════════════════════════════════════════════════════════
\begin{columns}
\column{0.42}

{
\colorlet{blocktitlebgcolor}{measureblue}
\block{Measurement Validity --- Does the metric measure what it claims?}{
  \renewcommand{\arraystretch}{1.25}
  {\small
  \begin{tabular}{@{}c l p{0.38\linewidth} l@{}}
    \toprule
    \textbf{ID} & \textbf{Name} & \textbf{Tests} & \textbf{Source} \\
    \midrule
    M1 & Reliability & Repeated measures agree? & Psychometrics \\
    M2 & Baseline sep.\ & Distinguishable from random? & Psychometrics \\
    M3 & Stability & Robust to perturbation? & Psychometrics \\
    M4 & Calibration & Numbers meaningful absolutely? & Psychometrics \\
    M5 & Sensitivity & Detects known-true effects? & Psychometrics \\
    M6 & Invariance & Consistent across conditions? & Psychometrics \\
    M7 & Selection corr.\ & Multiplicity controlled? & Statistics \\
    \bottomrule
  \end{tabular}}
}
}

\column{0.29}

{
\colorlet{blocktitlebgcolor}{measureblue}
\block{M2: Three Baseline Failures}{
  \textbf{SAE interpretability}: same distribution on
  random, untrained models as on trained ones. The
  metric cannot distinguish features from noise.

  \vspace{3mm}
  \textbf{IIA baselines} (Sutter): interchange
  intervention accuracy exceeds majority-class rates
  before any causal structure is identified.

  \vspace{3mm}
  \textbf{SAE reconstruction} (Heap): scores
  indistinguishable from a baseline ignoring the
  dictionary.
}
}

\column{0.29}

{
\colorlet{blocktitlebgcolor}{measureblue}
\block{M6: IOI Faithfulness}{
  The same circuit, same concept (faithfulness),
  produces scores spanning \textbf{below 0\% to over
  100\%} across six methodological choices.

  \vspace{3mm}
  If the reading changes when you swap the instrument,
  the reading is a property of the instrument, not the
  circuit.

  \vspace{3mm}
  \textbf{M4} (calibration) is attempted in every
  audited claim and confirmed in none.
}
}

\end{columns}

% ════════════════════════════════════════════════════════
% ROW 3: INTERNAL VALIDITY
% ════════════════════════════════════════════════════════
\begin{columns}
\column{0.42}

{
\colorlet{blocktitlebgcolor}{internalorange}
\block{Internal Validity --- Is the causal claim warranted?}{
  \renewcommand{\arraystretch}{1.15}
  {\small
  \begin{tabular}{@{}c l p{0.33\linewidth} l@{}}
    \toprule
    \textbf{ID} & \textbf{Name} & \textbf{Tests} & \textbf{Source} \\
    \midrule
    I1 & Necessity & Is the circuit required? & Neuroscience \\
    I2 & Sufficiency & Is the circuit enough? & Neuroscience \\
    I3 & Minimality & Every component earns its place? & Causal inf.\ \\
    I4 & Specificity & More effect than matched controls? & Medicine \\
    I5 & Rival exclusion & \emph{The} mechanism or \emph{a} mechanism? & Phil.\ of sci.\ \\
    I6 & Double dissoc.\ & Two interventions, crossed effects? & Neuroscience \\
    I7 & Confound control & Alternatives ruled out? & Causal inf.\ \\
    I8 & Confound sensit.\ & Unmeasured confounder strength? & Causal inf.\ \\
    I9 & Epistasis & Non-additive interaction? & Genetics \\
    I10 & Rescue & Restoring component recovers behav.? & Genetics \\
    I11 & Onset coupling & Mechanism appears with capability? & Med.\ micro.\ \\
    I12 & Offset coupling & Mechanism goes when cap.\ removed? & Med.\ micro.\ \\
    \bottomrule
  \end{tabular}}
}
}

\column{0.29}

{
\colorlet{blocktitlebgcolor}{internalorange}
\block{I4: Specificity / ORBITA}{
  The ORBITA trial (2017): cardiac stenting for stable
  angina produced no benefit over sham surgery. The
  intervention was real; specificity had never been tested.

  \vspace{3mm}
  \textbf{Knowledge Neurons}: off-target effects
  \textbf{worse than random}. The circuit is less specific
  than chance.

  \vspace{3mm}
  \textbf{IOI}: specificity against a size-matched random
  subnetwork has never been tested.
}
}

\column{0.29}

{
\colorlet{blocktitlebgcolor}{internalorange}
\block{I6: Double Dissociation}{
  \renewcommand{\arraystretch}{1.15}
  {\small
  \begin{tabular}{@{}lcc@{}}
    \toprule
    & \textbf{Behav.\ X} & \textbf{Behav.\ Y} \\
    \midrule
    Ablate A & Large & Small \\
    Ablate B & Small & Large \\
    \bottomrule
  \end{tabular}}

  \vspace{3mm}
  One row is a single dissociation. The crossed pattern
  rules out generalized disruption.

  \vspace{3mm}
  Only induction heads (token copying) meet this
  criterion, achieved by Feucht et al.\ (2025)---three
  years after the original paper.
}
}

\end{columns}

% ════════════════════════════════════════════════════════
% ROW 4: EXTERNAL VALIDITY
% ════════════════════════════════════════════════════════
\begin{columns}
\column{0.42}

{
\colorlet{blocktitlebgcolor}{externalpurple}
\block{External Validity --- Does the finding generalize?}{
  \renewcommand{\arraystretch}{1.25}
  {\small
  \begin{tabular}{@{}c l p{0.38\linewidth} l@{}}
    \toprule
    \textbf{ID} & \textbf{Name} & \textbf{Tests} & \textbf{Source} \\
    \midrule
    E1 & Intervention reach & Do two methods agree? & Causal inf.\ \\
    E2 & Prompt general.\ & Works on diverse prompts? & Causal inf.\ \\
    E3 & Cross-task & Transfers to related tasks? & Causal inf.\ \\
    E4 & Cross-model & Appears in other models? & Causal inf.\ \\
    E5 & Graded response & Partial ablation, partial effect? & Pharmacology \\
    E6 & Novel prediction & Predicts untested behaviors? & Phil.\ of sci.\ \\
    \bottomrule
  \end{tabular}}
}
}

\column{0.29}

{
\colorlet{blocktitlebgcolor}{externalpurple}
\block{E1: IOI Method Disagreement}{
  Six evaluations of the same circuit produce
  faithfulness scores spanning \textbf{below 0\%
  to over 100\%}.

  \vspace{3mm}
  When independent methods disagree on the
  \emph{direction} of an effect, the result is a
  property of the method, not the circuit.

  \vspace{3mm}
  The field's most studied mechanism shows the widest
  disagreement.
}
}

\column{0.29}

{
\colorlet{blocktitlebgcolor}{externalpurple}
\block{E5: Dose--Response}{
  The curve between zero and complete ablation
  distinguishes thresholds from graded mechanisms,
  reveals compensation and saturation, and detects
  non-monotonicity.

  \vspace{3mm}
  Most evaluations report a single ablation point.
  The interpolation is available in standard frameworks
  but almost never recorded.
}
}

\end{columns}

% ════════════════════════════════════════════════════════
% ROW 5: INTERPRETIVE VALIDITY + TAKEAWAY
% ════════════════════════════════════════════════════════
\begin{columns}
\column{0.42}

{
\colorlet{blocktitlebgcolor}{interpteal}
\block{Interpretive Validity --- Is the interpretation justified?}{
  {\small\textbf{Novel.} No existing validity framework
  includes these criteria.}

  \vspace{4mm}
  \renewcommand{\arraystretch}{1.25}
  {\small
  \begin{tabular}{@{}c l p{0.38\linewidth} l@{}}
    \toprule
    \textbf{ID} & \textbf{Name} & \textbf{Tests} & \textbf{Source} \\
    \midrule
    V1 & Level declaration & At what level is the claim? & Marr (1982) \\
    V2 & Level-evidence & Evidence supports that level? & Messick (1995) \\
    V3 & Alternative level & Explained at a different level? & Marr (1982) \\
    V4 & Unlicensed labeling & Name imports unmeasured property? & Messick (1995) \\
    V5 & Scope declaration & What does the claim not cover? & Reporting std.\ \\
    \bottomrule
  \end{tabular}}
}
}

\column{0.29}

{
\colorlet{blocktitlebgcolor}{interpteal}
\block{V4: Unlicensed Labeling}{
  \textbf{Othello ``world model''}: a probe recovers board
  state. The label imports spatial reasoning and planning
  that were never contrastively tested. One checkpoint
  passes; the other fails.

  \vspace{3mm}
  \textbf{Steering ``refusal direction''}: controls refusal
  across 13 models. But a lever is not a representation.
  The label imports content the evidence does not establish.
}
}

\column{0.29}

{
\colorlet{blocktitlebgcolor}{colorOne}
\block{Takeaway}{
  {\large\bfseries A mechanistic claim is only as strong
  as its weakest validity dimension.}

  \vspace{3mm}
  \begin{itemize}[leftmargin=*, itemsep=2mm]
    \item \textbf{36} criteria, \textbf{5} validity types,
      \textbf{8} source disciplines
    \item \textbf{16} claims from \textbf{15} papers audited
    \item The field's most studied circuit (IOI) is
      disconfirmed on stability, invariance, and prompt
      generalization
    \item Interpretive validity (V1--V5) has no counterpart
      in existing frameworks
  \end{itemize}

  \vspace{3mm}
  \begin{center}
  \texttt{elliot@elliottower.ai}

  \vspace{2mm}
  {\small Paper:}
  \texttt{doi.org/10.5281/zenodo.20478480}
  \end{center}
}
}

\end{columns}

\end{document}
